\section{Zadanie 6}

Niech \(X\) będzie zmienną losową liczącą liczbę kopii kliki \(K_4\) na czterech wierzchołkach w grafie wylosowanym z modelu \(G_{n,p}\) (to znaczy, mając \(n \geq 4\) wierzchołków dla każdej pary wierzchołków losujemy niezależnie z prawdopodobieństwem \(p \in (0, 1)\), czy są połączone krawędzią, czy nie). Oblicz wartość oczekiwaną i wariancję \(X\).

\subsection*{Rozwiązanie}

Wartość oczekiwana jest prosta do wyliczenia. Wystarczy zdefiniować indykatory \(\{X_I\}_{I \in \binom{V}{4}}\) dla każdego 4-elementowego podzbioru wierzchołków, które są równe 1 wpp. gdy ten zbiór indukuje kopię \(K_4\). Otrzymujemy wartość oczekiwaną:

\[
E[X] = E[\sum_{I \in \binom{V}{4}}X_I] = \sum_{I \in \binom V 4}E[X_I] = \sum_{I \in \binom V 4}p^6 = \binom{n}{4}p^6
\]

Z wariancją jest nieco więcej zabawy, ale na szczęście możemy zauważyć, że indykatory zbiorów o 0 lub 1 wspólnych wierzchołkach byłyby rozłączne krawędziowo, więc byłyby niezależne (a więc \(Cov(X_I, X_J) = 0\)). Możemy więc wyliczyć wariancję pierwsze poznając wartości kowariancji dla przecinających się krawędziowo zbiorów.

\subsubsection*{2 wspólne wierzchołki}

Niech \(X, Y \in \binom{V}{4}, |X \cap Y| = 2\).

\begin{align*}
    Cov(X_I, X_J) &= E[X_IX_J] - E[X_I]E[X_J] =\\
    &= p^{11} - (p^{6})^2 = (1-p)p^{11}
\end{align*}

\subsubsection*{3 wspólne wierzchołki}

Niech \(X, Y \in \binom{V}{4}, |X \cap Y| = 3\)

\begin{align*}
    Cov(X_I, X_J) &= E[X_IX_J] - E[X_I]E[X_J] =\\
    &= p^9 - (p^6)^2 = (1-p^3)p^9
\end{align*}

\subsubsection*{Wariancje}

Pozostaje jeszcze tylko wyliczyć wariancję pojedynczego indykatora:

\[
Var(X_I) = E[X_I^2] - E[X_I]^2 = p^6 - p^{12} = (1-p^6)p^6
\]

I ostateczny wynik:

\begin{align*}
Var(X) &= Var(\sum_{I \in \binom V 4}X_I) = \sum_{I \in \binom V 4}Var(X_I) + \sum_{I, J \in \binom V 4, I \neq J}Cov(X_I, X_J)\\
&= \sum_{I \in \binom V 4}Var(X_I) + \sum_{I, J \in \binom V 4, I \cap J = 2}Cov(X_I, X_J) + \sum_{I, J \in \binom V 4, I \cap J = 3}Cov(X_I, X_J)\\
&= \sum_{I \in \binom V 4}Var(X_I) + \sum_{I, J \in \binom V 4, I \cap J = 2}Cov(X_I, X_J) + \sum_{I, J \in \binom V 4, I \cap J = 3}Cov(X_I, X_J)\\
&= \binom{n}{4}(1-p^6)p^6 + \binom{n}{6}\binom{6}{2}\binom{4}{2}(1-p)p^{11} + \binom{n}{5}\binom{5}{3}\binom{2}{1}(1-p^3)p^9\\
&= \binom{n}{4}(1-p^6)p^6 + 90\binom{n}{6}(1-p)p^{11} + 20\binom{n}{5}(1-p^3)p^9\\
\end{align*}

\subsection*{Odpowiedź}
\[
\boxed{E[X] = \binom{n}{4}p^6, Var[X] \text{ powyżej}}
\]
